\section{Full preorderings and fractional-cardinality moments}
\label{sec:cardinality-preordering}
An order-$r$ oracle below acts on polynomials of total degree at most
$2r$, where $r\ge1$ is an integer. On a box
$B=\prod_i[a_i,b_i]\subseteq[0,1]^n$, it consists of linear functionals
$L:\R[x]_{\le2r}\to\R$ satisfying
\begin{align}
 L[1]&=1,\label{eq:preorder-normalization}\\
 L[e v]&=0\quad(\deg v\le2r-1),\qquad
 e=\sum_i x_i-K,\label{eq:preorder-equality}\\
 L[g P^2]&\ge0\quad
 \left(\sum_{j=1}^a\deg g_j+2\deg P\le2r,\quad
 g=\prod_{j=1}^a g_j\right),\label{eq:preorder-positivity}
\end{align}
where each $g_j$ is a coordinate-bound slack $x_i-a_i$ or $b_i-x_i$.
The empty product and repeated slack factors are included, and the square
polynomial $P$ may couple every coordinate. Zero polynomials impose no
condition; constants have degree zero. The node bound
$\operatorname{LB}_r(B)$ is the infimum of $L[F]$ over these functionals.
This is the full truncated box preordering, with all allowed products
of slacks, not only single-generator localizers. Evaluation at a feasible
point proves validity.

At $r=1$ this model is exactly
\eqref{eq:sdp-model}--\eqref{eq:box-rlt}: positivity on affine squares
is augmented PSD, products of at most two slacks give box RLT, single
slacks give first-moment bounds, and the balance times affine polynomials
gives its first and second moment equations. At every order $r\ge1$ the
model dominates that degree-two model and hence the chord oracle.
Certification again means an infeasible region or
$\operatorname{LB}_r(B)\ge T_*=1/4-\epsilon$ for
problem~\eqref{eq:cardinality-problem}. Original dimension remains $n$.

\subsection{A self-contained positivity proof}
The moments below are the classical fractional-cardinality moments
associated with Grigoriev's knapsack lower bound
\citep{Grigoriev2001}. Potechin gives the explicit construction and its
classical provenance \citep[Theorem 1, Example 18, Theorem 44]{Potechin2019}.
We prove positivity directly in the sufficient range needed for the full
preordering, without an external positivity theorem.

For real $v$ and integer $a\ge0$, write
$(v)_a=v(v-1)\cdots(v-a+1)$, with $(v)_0=1$.
In $s$ formal Boolean variables $u_1,\ldots,u_s$, let $\operatorname{red}$
replace every positive power of $u_i$ by $u_i$. This is the canonical
multilinear representative modulo $u_i^2-u_i$; it respects sums and
products modulo these identities and cannot increase degree. For
$u_S=\prod_{i\in S}u_i$ define
\begin{equation}\label{eq:fractional-moments}
 E_{s,t}[u_S]=\frac{(t)_{|S|}}{(s)_{|S|}},\qquad
 E_{s,t}[P]=E_{s,t}[\operatorname{red}P].
\end{equation}
All subset sizes are at most $s$, so denominators are positive. The
functional is defined on the Boolean quotient and is used only through
the degree stated in each assertion; positivity on all polynomials is not
assumed. In particular $E_{s,t}[1]=1$. For $s=0$, deterministic evaluation
of constants will be used instead of a fractional-cardinality functional.

\begin{lemma}[Cardinality identities and moment positivity]
\label{lem:fractional-positive}
Let $d\ge1$ be an integer, $s\ge2d$ an integer, and
$t,s-t\ge2d-1$. Then
\begin{equation}\label{eq:fractional-cardinality-identity}
 E_{s,t}\left[\left(\sum_i u_i-t\right)v\right]=0
 \quad(\deg v\le2d-1),
\end{equation}
and $E_{s,t}[P^2]\ge0$ for every $\deg P\le d$.
\end{lemma}
\begin{proof}
For squarefree $u_S$ of degree $a\le2d-1\le s-1$,
\[
 E_{s,t}\left[\left(\sum_i u_i\right)u_S\right]
 =a\frac{(t)_a}{(s)_a}
  +(s-a)\frac{(t)_{a+1}}{(s)_{a+1}}
 =t\frac{(t)_a}{(s)_a}.
\]
Boolean reduction and linearity prove
\eqref{eq:fractional-cardinality-identity} for every allowed multiplier.

To reduce positivity to homogeneous polynomials, for $|S|=a\le d$ put
\begin{equation}\label{eq:fractional-homogenization}
 H_S(u)=\frac{\sum_{T\supseteq S,\ |T|=d}u_T}
                  {\binom{t-a}{d-a}}.
\end{equation}
The denominator is positive under $t\ge2d-1$ (and is one for $a=d$).
In the Boolean quotient the numerator equals
$u_S\binom{\sum_i u_i-a}{d-a}$: at a Boolean vector it counts the
$d$-subsets containing $S$ when all coordinates of $S$ equal one, and
both sides vanish otherwise. The polynomial
$\binom{z-a}{d-a}-\binom{t-a}{d-a}$ is divisible by $z-t$ with quotient
degree at most $d-a-1$ when $a<d$. Consequently
\[
 u_S-H_S=\left(\sum_i u_i-t\right)Q_S
 \quad\hbox{in the Boolean quotient},\qquad \deg Q_S\le d-1.
\]
Replace each monomial of $\operatorname{red}P$ by its $H_S$, obtaining a
homogeneous squarefree polynomial $H$ of degree $d$. Then
$P-H=(\sum_i u_i-t)Q$ modulo Boolean identities with $\deg Q\le d-1$.
Since $\deg(Q(P+H))\le2d-1$, the cardinality identity gives
$E_{s,t}[P^2]=E_{s,t}[H^2]$. Every expression used here has degree at most
$2d$ before Boolean reduction.

Index the homogeneous moment matrix $N$ by $d$-subsets $I,J$ of $[s]$:
\[
 N(I,J)=\frac{(t)_{2d-|I\cap J|}}{(s)_{2d-|I\cap J|}}.
\]
Let $P_j(I,J)=\binom{|I\cap J|}{j}$ for $0\le j\le d$.
This is a Gram matrix: associate with $I$ the vector indexed by
$j$-subsets $A$, with entry $\ind\{A\subseteq I\}$.
The exact decomposition is
\begin{equation}\label{eq:fractional-gram}
 N=\sum_{j=0}^d
   \frac{(t)_{2d-j}(s-t)_j}{(s)_{2d}}P_j.
\end{equation}
Indeed, in an entry with $|I\cap J|=\ell$, the numerator on the right
is
\begin{align*}
 \sum_{j=0}^\ell\binom\ell j(t)_{2d-j}(s-t)_j
 &=(t)_{2d-\ell}
   \sum_{j=0}^\ell\binom\ell j
       (t-2d+\ell)_{\ell-j}(s-t)_j\\
 &=(t)_{2d-\ell}(s-2d+\ell)_\ell.
\end{align*}
The last equality is the falling-factorial Vandermonde identity: count
ordered choices from two disjoint finite sets for nonnegative integer
arguments, then extend to real arguments as a polynomial identity.
Division by $(s)_{2d}$ gives the claimed entry. No division by
$(t)_{2d-\ell}$ was used, so zero coefficients cause no exception.
All coefficients in \eqref{eq:fractional-gram} are nonnegative because
$t,s-t\ge2d-1$, and $(s)_{2d}>0$. Thus $N\succeq0$ and
$E_{s,t}[H^2]\ge0$.
\end{proof}

\begin{lemma}[Assignment indicators and every slack product]
\label{lem:assignment-localizer}
Suppose $r\ge1$, $s\ge2r$, and $t,s-t\ge2r-1$.
For every product $g$ of $u_i$ and $1-u_i$, including repetitions, and
every polynomial $P$ with
$\deg g+2\deg P\le2r$, one has $E_{s,t}[gP^2]\ge0$.
\end{lemma}
\begin{proof}
Boolean reduction makes $g$ zero if both $u_i$ and $1-u_i$ occur for
some $i$. Otherwise it becomes an assignment indicator
\[
 I_{A,B}=\prod_{i\in A}u_i\prod_{j\in B}(1-u_j),\qquad A\cap B=\varnothing.
\]
Put $a=|A|$, $b=|B|$, and $v=a+b\le\deg g$. For a squarefree monomial
$u_C$ on the remaining coordinates, direct expansion gives
\begin{equation}\label{eq:indicator-moment}
 E_{s,t}[I_{A,B}u_C]
 =\frac{(t)_{a+|C|}(s-t)_b}{(s)_{v+|C|}}
 \quad(v+|C|\le s).
\end{equation}
One verification is induction on $b$: inserting an additional factor
$1-u_j$ subtracts the expression with $j$ also fixed to one, and putting
the two terms over their common denominator supplies the next factor
$s-t-b$. The case $b=0$ is \eqref{eq:fractional-moments}.
In particular the indicator weight is
\[
 \pi=E_{s,t}[I_{A,B}]=\frac{(t)_a(s-t)_b}{(s)_v}\ge0.
\]
If $P$ is constant, the desired inequality is $\pi P^2\ge0$.
If $P$ is nonconstant, $v\le\deg g\le2r-2$, so all factors in this
weight are strictly positive. Write
$\widehat P=P(1_A,0_B,u_{[s]\setminus(A\cup B)})$.
For $v<s$, cancellation in \eqref{eq:indicator-moment} proves
\begin{equation}\label{eq:conditional-moments}
 E_{s,t}[I_{A,B}P^2]=\pi E_{s-v,t-a}[\widehat P^2].
\end{equation}
This identity is used only when its remaining moments are defined.
If $\widehat P$ is constant, the value is simply $\pi\widehat P^2$;
the same direct evaluation applies when $v=s$, so no zero-variable
conditional functional is needed.
Otherwise let $d=\deg\widehat P\ge1$. Then
\[
 v+2d\le2r,\quad s-v\ge2d,\quad
 t-a\ge2r-1-a\ge2d-1,\quad s-t-b\ge2d-1.
\]
Lemma~\ref{lem:fractional-positive} applies to
$E_{s-v,t-a}$. This proves positivity, including for arbitrary squares
that couple all unassigned variables and for repeated original slacks.
\end{proof}

The sufficient degree range in Lemma~\ref{lem:fractional-positive} is
not the sharp classical range for the moment matrix alone. For the
\emph{full preordering} and this functional, the slack products explain
the two-sided restriction: if $0<t<s$ is noninteger,
$s\ge2r$, and $t<2r-1$, a product of $\lfloor t\rfloor+2\le2r$
distinct lower slacks has a negative falling-factorial moment. Applying
the same argument to upper slacks treats $s-t<2r-1$. Integer $t$ is
different: \eqref{eq:fractional-moments} is then expectation under the
uniform distribution of $t$-subsets. Sharper moment-matrix positivity
alone therefore does not extend this full-preordering construction past
the noninteger slack-product obstruction.

\subsection{The order versus region tradeoff}
\begin{theorem}[Full-preordering spatial cover bound]
\label{thm:preordering-cover}
For problem~\eqref{eq:cardinality-problem}, let
$n=k+m+z$ with positive integers $k,m,z$, let $r\ge1$ be an integer,
and let $0\le\epsilon<1/4$. Put
\begin{equation}\label{eq:qr}
 q_r=\min\{k-2r+2,\ z-2r+2,\ m(1/2-2\epsilon)\}.
\end{equation}
If $q_r>0$, every box cover certified by
\eqref{eq:preorder-normalization}--\eqref{eq:preorder-positivity} at
$T_*=1/4-\epsilon$ has at least $\mathcal N(q_r)$ members.
If $q_r\le0$, only the trivial bound of one region is asserted.
\end{theorem}
\begin{proof}
Take a box containing a witness~\eqref{eq:cardinality-witness} and
suppose $|R|<q_r$. Since $|R|$ and $k-2r+2$ are integers,
$|R|<k-2r+2$ implies $|H\cap U|\ge k-|R|\ge2r-1$;
likewise $|Z\cap U|\ge2r-1$. Set
\[
 s=|U|,\qquad t=|H\cap U|+p|M\cap U|.
\]
Then $t,s-t\ge2r-1$. Also $s\ge2r$: for $r=1$ a unit and a zero
coordinate remain, and for $r\ge2$ we have $s\ge4r-2\ge2r$.
Define $L$ by fixing $x_R=w_R$ and applying $E_{s,t}$ to $x_U=u$.
This degree-preserving substitution is normalized. It changes $e$ into
$\sum_{i\in U}u_i-t$, whose products with every substituted multiplier
of degree at most $2r-1$ vanish by
Lemma~\ref{lem:fractional-positive}. Bound slacks on $R$ are
nonnegative constants, and those on $U$ are Boolean slacks. Removing
constant factors cannot increase degrees, so
Lemma~\ref{lem:assignment-localizer} proves every required positivity
constraint.

The Boolean penalty on $U$ vanishes, leaving
\begin{equation}\label{eq:restricted-penalty}
 L[F]=|M\cap R|p(1-p)\le |R|p(1-p)
 <\frac{q_r}{2m}\le\tfrac14-\epsilon.
\end{equation}
The strict inequality follows from $|R|<q_r$ and $q_r>0$, including
$R=\varnothing$. This feasible functional precludes certification.
Thus every certified box containing a witness has $|R|\ge q_r$,
and Lemma~\ref{lem:endpoint-count} proves the bound. Non-strict pruning
at the target is covered because the constructed value is strictly below
it, and no integer rounding of $q_r/2$ is needed.
\end{proof}

For balanced parts $k=m=z=t$, positivity of $q_r$ is precisely
$r<(t+2)/2$ together with $\epsilon<1/4$. A uniform exponential
bound as $n=3t\to\infty$ follows whenever $\epsilon<1/4$ is fixed
and $r\le ct$ for a fixed $c<1/2$. At $\epsilon=1/8$, the exponent
$n/24$ survives whenever
\[
 t-2r+2\ge t/4,\qquad\hbox{equivalently}\qquad r\le3t/8+1.
\]
An order near the positivity boundary can give only a constant $q_r$,
so positivity of $q_r$ alone does not give exponential size.
When $\epsilon=1/4$, the root is certified because the mixed slacks give
$L[F]\ge0$. None of these cases divides by $T_*$ or by a vanishing $q_r$.
Increasing $r$ may make the oracle itself costly; the theorem grants the
oracle and counts regions only.

The affine-inequality closure following Theorem~\ref{thm:sdp-cover}
extends to this degree. Expand a product of valid affine inequalities
in their nonnegative box-slack and balance representations and multiply
by an allowed square. Terms with a balance factor vanish within degree
$2r$, and all remaining terms are existing preordering inequalities.
This still grants no arbitrary coupled nonlinear valid cut, no localizer
of an objective cutoff, and no analytic solution of the coupled problem.
The charged-cover convention of Proposition~\ref{prop:charged-tolerance}
applies to objective-based spatial reductions.

\subsection{Unique minimizers and increasing concave costs}
\begin{corollary}[Distinct perturbations]\label{cor:unique-cardinality}
Let $0<d_1<\cdots<d_n$ and $\sum_i d_i\le\eta$, where $\eta>0$.
On $\mathcal F_{n,k}$ replace $F$ by
$F_d=F+\sum_i d_ix_i$. Its unique optimizer has
\begin{equation}\label{eq:unique-cardinality}
 x_i^*=1\ (i\le k),\qquad x_{k+1}^*=1/2,\qquad
 x_i^*=0\ (i\ge k+2),
\end{equation}
with value $F_d^*=1/4+\sum_{i\le k}d_i+d_{k+1}/2$.
For the cover assertion take $n=k+m+z$ with positive integers $k,m,z$
and integer $r\ge1$, as in Theorem~\ref{thm:preordering-cover}.
Under the order-$r$ box oracle, a cover at
$T_*=F_d^*-\epsilon$ has at least $\mathcal N(q_{r,\eta})$ members
provided $0\le\epsilon$, $\epsilon+\eta<1/4$, and
\begin{equation}\label{eq:perturbed-qr}
 q_{r,\eta}=\min\{k-2r+2,\ z-2r+2,
                         m(1/2-2\epsilon-2\eta)\}>0.
\end{equation}
There is no nonidentity coordinate permutation preserving both feasible
set and objective, even when objective equality is required only on the
demand slice.
\end{corollary}
\begin{proof}
The Hessian of $F_d$ is $-2I$, so a nonvertex, lying strictly between
two distinct feasible points, cannot minimize this strictly concave
function. All vertices have base value $1/4$. Among their linear costs,
an exchange puts the largest coordinate value at the smallest coefficient:
if $i<j$ and $x_i<x_j$, exchanging them reduces cost by
$(d_j-d_i)(x_j-x_i)>0$. Hence the unique minimizing vertex is
\eqref{eq:unique-cardinality}. It also uniquely minimizes the linear
cost over the whole slice: if $i<j$, $x_i<1$, and $x_j>0$, transferring
$\min\{1-x_i,x_j\}>0$ from coordinate $j$ to coordinate $i$ strictly
decreases that cost. At a linear-cost minimizer no such pair remains,
so the balance forces exactly \eqref{eq:unique-cardinality}.

The node constraints themselves impose $0\le L[x_i]\le1$, since
$x_i$ and $1-x_i$ are nonnegative combinations of node slacks and
constants. The constructed functional therefore gives linear cost at
most $\eta$. If $|R|<q_{r,\eta}$, its total value is strictly below
\[
 \frac{m(1/2-2\epsilon-2\eta)}{2m}+\eta
 =1/4-\epsilon\le F_d^*-\epsilon,
\]
so the previous count applies.

The slice has an interior point relative to its balance hyperplane,
namely $x_i=K/n\in(0,1)$. A linear form constant on this slice has
coefficient vector proportional to the all-ones vector. If a coordinate
permutation preserves the objective there, the permuted coefficient
vector differs from $d$ by a constant vector. Comparing sums forces this
constant to be zero, and distinctness forces the identity permutation.
\end{proof}

An explicit choice is $d_i=2\eta i/[n(n+1)]$, whose sum is $\eta$.
For $n=3t$, $\epsilon=1/8$, and $\eta=1/32$, the lower bound is
$(3/2)^{n/32}$ whenever $r\le13t/32+1$.
If $\eta<\epsilon$, the chord certificate at tolerance
$\epsilon-\eta$ has base bound at least
$1/4+\eta-\epsilon\ge F_d^*-\epsilon$; the nonnegative linear term
preserves it. Thus this explicit perturbed family also has
$2^{\Theta(n)}$ certificate size under the specified oracles.

Adding the conserved flow gives increasing strictly concave coordinate
costs:
\begin{equation}\label{eq:increasing-allocation}
 C_d(x)=\sum_i[x_i(2-x_i)+d_ix_i]=F_d(x)+K
 \quad(x\in\mathcal F_{n,k}).
\end{equation}
Each derivative is $2-2x_i+d_i>0$ on $[0,1]$ and each second derivative
is $-2$. Every equality-preserving functional adds exactly $K$, so the
optimizer and all absolute-gap conclusions are unchanged. When $k$ is
proportional to $n$, this raises the optimum to order $n$, while the
absolute gap in the obstruction stays constant. A fixed relative tolerance
therefore does not follow from this absolute bound;
Section~\ref{sec:relative-blocks} treats relative tolerance with
fixed-dimensional blocks.
